\section{Introduction}
\label{sec:intro}

Paragraph 1: Introduce the broader research problem and explain why minority opinions matter  What is minority opinion? Why is minority opinion important? For example: A functioning democratic dialogue requires dissenting opinions to be present, heard, and considered.

Paragraph 2: Explain how this research question has been studied, what are the limitations in laboratory environments? motivate the use of social media

Paragraph 3: Introduce AITA and existing research on minority opinions in online communities. What is AITA,  What gaps remain in the existing literature? 

Paragraph 4: Define our focus, research questions, and approach

Core research question: How are minority opinions expressed and treated in online moral discussions? 

\begin{enumerate}
    \item \textbf{RQ1:} Does the difference in how majority and minority
    opinions are treated vary between controversial and
    non-controversial posts?

    \item \textbf{RQ2:} How do users respond to minority opinions in
    non-controversial posts?

    \item \textbf{RQ3:} Does the treatment of minority opinions depend
    on whether they accuse or defend the post author?

    \item \textbf{RQ4:} Do minority opinions develop differently in
    AITA and AITAH?
\end{enumerate}

Paragraph 5: Findings and Contributions




Online platforms make minority opinions visible, but visibility does not necessarily mean equal treatment. A dissenting comment may be rejected through downvotes while simultaneously attracting substantial attention through replies. Conversely, a widely supported comment may receive little direct engagement because it leaves little to contest. Understanding how crowds treat minority voices therefore requires separating approval from conversational attention.
Not every numerically smaller opinion, however, constitutes meaningful dissent. When opinion is evenly divided, the smaller side is simply one participant in an unresolved debate. When more than 80\% of a crowd supports one judgment, expressing the opposite view becomes a challenge to an established consensus. Minority status is thus relational: the same judgment can be an ordinary competing opinion in a divided discussion and consequential dissent in a strong-consensus discussion.
We study this distinction in two Reddit communities devoted to interpersonal moral judgment, r/AmItheAsshole (AITA) and r/AITAH. The comparison also provides an institutional contrast. AITA formalizes the prevailing judgment as an official verdict flair approximately 18 hours after posting, whereas AITAH has no comparable rule. Together, these settings allow us to examine when minority judgments become consequential, how crowds respond to them, and how dissent changes as consensus develops over time.

\subsection{Research Questions}

\textbf{RQ1. Non-controversial vs Controversial: When does a minority judgment become consequential dissent?}
Does the treatment of minority comments depend on the strength of crowd consensus? Specifically, are majority–minority differences larger in non-controversial posts than in controversial posts?

\textbf{RQ2. Reply/Challenge behavior: How do crowds treat dissent against strong consensus?}
Are minority judgments sanctioned through voting while being amplified through replies and deeper discussion? Does dissent primarily initiate conversation, or does it also intensify exchanges once they have begun?

\textbf{RQ3. Accusing/Defending: Does the moral direction of dissent matter?}
Are comments that accuse someone whom the crowd has largely absolved treated differently from comments that defend someone whom the crowd has largely condemned?

\textbf{RQ4. AITA/AITAH: How does dissent change as consensus develops and becomes formalized?}
How do the prevalence and expressed certainty of minority judgments evolve over time? Do these dynamics differ between AITA, where a verdict is formally displayed after 18 hours, and AITAH, where no official verdict is imposed?

\subsection{Method}
We analyze more than 15 million top-level judgment comments: over 11 million from AITA and over 4 million from AITAH. Judgments are grouped into two opposing categories: YA, indicating that the focal actor is at fault, and NA, indicating that the actor is not at fault.

Posts are classified by their YA share:
Non-controversial -- YA share below 20\% or above 80\% versus Controversial -- YA share between 40\% and 60\%.
Within each post, a comment is classified as majority or minority according to whether its judgment agrees with the dominant side. In non-controversial posts, minority comments represent dissent against strong consensus. In controversial posts, the minority label identifies the numerically smaller side of a divided discussion and serves as a comparison condition.

We examine four dimensions of treatment:
voting outcomes, measured by comment score;
conversational response, measured by reply probability, reply count, and reply-tree depth;
conversational value, measured by the aggregate score of the reply subtree;
expressed certainty and temporal change, using certainty annotations from approximately 1,000 large threads.
We also distinguish accusatory dissent—YA judgments in NA-dominant posts—from defensive dissent—NA judgments in YA-dominant posts.

\subsection{Main Findings}
First, minority status becomes consequential primarily under strong consensus. In non-controversial posts, minority comments receive substantially more replies than majority comments. In controversial posts, however, majority–minority differences are small or absent. The crowd therefore responds not to numerical minority status alone, but to an opinion’s position against a clear collective judgment.

Second, strong-consensus dissent is sanctioned through voting but amplified through conversation. Minority comments receive much lower scores, including negative average scores in AITA, yet are three to five times more likely to receive replies. They are also approximately six times more likely to initiate deeper exchanges. Most of this difference occurs at conversational entry: minority comments are especially likely to start a discussion, but majority–minority differences become much smaller once an exchange is already underway.

Third, the engagement created by dissent does not necessarily reward the dissenter. Even when a minority comment has a negative score, the reply subtree rooted in that comment can accumulate a positive total score. The crowd may therefore punish the initiating judgment while positively evaluating the rebuttals and discussion it generates.

Fourth, dissent is asymmetric. Accusatory dissent is punished more heavily but generates roughly 1.7–1.9 times more discussion than defensive dissent. Accusations that challenge a lenient consensus provoke confrontation, whereas defenses that challenge a condemnatory consensus are more likely to be ignored.

Finally, dissent changes over time, but differently across the two communities. In AITA, minority certainty increases approximately linearly as discussions age, even though greater certainty does not produce greater voting rewards. AITAH shows no comparable linear increase. Importantly, AITA exhibits no visible certainty jump at the 18-hour verdict threshold; the post-verdict regression coefficient is statistically significant but substantively very small. The temporal pattern is therefore more consistent with gradual changes in the conversational environment or speaker composition than with an immediate verdict effect.

Taken together, the findings show that a minority opinion becomes consequential dissent when there is a strong crowd consensus to oppose. Once it does, the crowd rejects the judgment but engages the challenge—extracting conversation from dissent while leaving much of its social cost with the speaker.